\section{Related Work}
\label{sec:related}

\subsection{Behavioral Finance and Social Relations}

Prospect theory describes reference dependence, loss aversion, and diminishing sensitivity~\cite{kahneman1979prospect,tversky1992advances}. Mental accounting and endogenous references explain how one price can be assessed against different benchmarks~\cite{thaler1985mental,koszegi2006model}. The disposition effect motivates tracking acquisition costs separately from current prices~\cite{shefrin1985disposition,odean1998reluctant}. Civitas connects these financial states to person-to-person relations through portfolio-exposure similarity. Acquisition cost, recent highs, peer consensus, and policy signals enter reference-dependent valuation.

Social psychology studies group affiliation, emotional contagion, and informational and normative influence~\cite{tajfel1979integrative,hatfield1993emotional,deutsch1955study}. Small-world connectivity and selective exposure motivate the contact network and access to information~\cite{watts1998collective,sunstein2001republic}. At the person-to-group scale, Civitas uses attention-bounded message selection and synchronous panic transitions to update local social states. Confidence controls the weight placed on peers. Investor archetypes specify heterogeneous roles independently of network position. The modeled group feedback arises through neighborhood interaction.

\subsection{Coordination and Market Institutions}

Strategic complementarities, expectations about others' valuations, and public-information games describe incentives to coordinate~\cite{cooper1988coordinating,keynes1936general,morris2002social}. Information-cascade models explain how earlier actions influence later decisions~\cite{bikhchandani1992theory,banerjee1992simple}. Noise-trader risk and limits to arbitrage explain why price deviations can persist~\cite{delong1990noise,shleifer1997limits}. These theories motivate peer-dependent updates in Civitas, which evolves its specified behavioral rules forward in time. Its group-to-society scale is the modeled market institution, where collective orders produce prices and trading opportunities that feed back into individual decisions.

Exchange design governs this aggregation. Zero-intelligence experiments isolate the contribution of budget constraints and double-auction rules~\cite{gode1993allocative}. Interacting chartist and fundamentalist models generate heavy tails and volatility clustering~\cite{lux1999scaling,lux2000volatility,raberto2001agent}. Microstructure theory links price formation to information asymmetry~\cite{kyle1985continuous,glosten1985bid}. ABIDES provides a discrete-event limit order book with configurable communication delays and programmable agents~\cite{abides2020}. Civitas uses discrete decision barriers followed by sequential order arrivals at a continuous double auction, with a common information boundary before matching.

\subsection{Language Models in Market Simulations}

Generative agents interpret observations and enact roles through language~\cite{park2023generative}. Financial applications include layered memory in FinMem, multimodal signals and tools in FinAgent, and role-based trading in StockAgent and TradingAgents~\cite{yu2023finmem,zhang2024finagent,li2024stockagent,xiao2024tradingagents}. EconAgent models household labor and consumption~\cite{li2024econagent}. TwinMarket couples belief--desire--intention investors and social information to a call auction~\cite{yang2025twinmarket}. MarketSim separates strategic LLM decisions from high-frequency execution in a continuous double auction~\cite{marketsim2026}. StockSim provides order-level LLM evaluation with two execution settings, latency, and slippage~\cite{papadakis2025stocksim}.

Table~\ref{tab:related_comparison} separates exchange design, social and behavioral mechanisms, decision organization, and evaluation focus. Civitas combines independently adjustable social and valuation rules with integer accounting and branches from a shared state. The comparison follows changes from interpersonal exposure and group feedback to order intent and execution under a common decision contract.

\begin{table*}[t]
  \centering\small
  \setlength{\tabcolsep}{4pt}
  \caption{\tr{金融市场仿真平台的设计与评估重点。各项概括所引文献报告的功能与研究重点。}{Design and evaluation focus of financial market simulators. Entries summarize features and study emphases reported in the cited work.}}
  \label{tab:related_comparison}
  \begin{tabularx}{\textwidth}{@{}>{\raggedright\arraybackslash}p{1.95cm} >{\raggedright\arraybackslash}p{2.65cm} >{\raggedright\arraybackslash}p{3.5cm} >{\raggedright\arraybackslash}p{3.3cm} >{\raggedright\arraybackslash}X@{}}
    \toprule
    \tr{平台}{Platform} & \tr{交易机制}{Exchange} & \tr{社会与行为机制}{Social / behavioral layer} & \tr{决策组织}{Decision organization} & \tr{评估与审计重点}{Evaluation and audit focus} \\
    \midrule
    \textbf{ABIDES}\newline\cite{abides2020} & \tr{离散事件限价订单簿；通信延迟}{Discrete-event limit order book; communication latency} & \tr{背景交易者与策略交易者}{Background and strategic traders} & \tr{可编程主体策略}{Programmable agent policies} & \tr{订单与消息记录；交易实验}{Order and message traces; trading experiments} \\
    \addlinespace
    \textbf{TwinMarket}\newline\cite{yang2025twinmarket} & \tr{日级集合竞价}{Daily call auction} & \tr{交易相似度社会信息流；谣言传播}{Trading-similarity social feed; rumor diffusion} & \tr{基于信念--愿望--意图的 LLM 投资者}{Belief--desire--intention LLM investors} & \tr{风格化事实与社会响应场景}{Stylized facts and social-response scenarios} \\
    \addlinespace
    \textbf{MarketSim}\newline\cite{marketsim2026} & \tr{纳秒精度连续双向拍卖}{Nanosecond-resolution continuous double auction} & \tr{异质投资者；新闻与政策信息}{Heterogeneous investors; news and policy information} & \tr{战略层 LLM 决策与快速执行层}{Strategic LLM decisions with a faster execution layer} & \tr{风格化事实与高频价格跟踪}{Stylized facts and high-frequency price tracking} \\
    \addlinespace
    \textbf{StockSim}\newline\cite{papadakis2025stocksim} & \tr{订单级执行；延迟与滑点}{Order-level execution; latency and slippage} & \tr{异质交易角色与多主体协调}{Heterogeneous trading roles and multi-agent coordination} & \tr{两种执行设置下的 LLM 交易决策}{LLM trading decisions in two execution settings} & \tr{序列交易决策评估}{Sequential trading decision evaluation} \\
    \addlinespace
    \textbf{Civitas Economica} & \tr{同步决策；连续双向拍卖；整数结算}{Synchronous decisions; continuous double auction; integer settlement} & \tr{暴露加权联系；参考点估值；恐慌与从众}{Exposure-weighted ties; reference valuation; panic and herding} & \tr{共用契约下的规则、启发式、LLM 与双速策略}{Shared contract for rules, heuristics, LLM, and dual-speed policies} & \tr{共同状态分支；账本与接口检查}{Shared-state branches; ledger and interface checks} \\
    \bottomrule
  \end{tabularx}
\end{table*}

